%-------------------------
% Vinicius Cubas Brand -- Currículo
% Baseado no template de Jake Gutierrez
%-------------------------

\documentclass[letterpaper,11pt]{article}

\usepackage{latexsym}
\usepackage[empty]{fullpage}
\usepackage{titlesec}
\usepackage[usenames,dvipsnames]{color}
\usepackage{enumitem}
\usepackage[hidelinks]{hyperref}
\usepackage{fancyhdr}
\usepackage[brazil]{babel}
\usepackage[utf8]{inputenc}
\usepackage[T1]{fontenc}
\usepackage{tabularx}
\input{glyphtounicode}

%----------CONFIGURAÇÃO DE PÁGINA----------
\pagestyle{fancy}
\fancyhf{}
\fancyfoot{}
\renewcommand{\headrulewidth}{0pt}
\renewcommand{\footrulewidth}{0pt}

\addtolength{\oddsidemargin}{-0.5in}
\addtolength{\evensidemargin}{-0.5in}
\addtolength{\textwidth}{1in}
\addtolength{\topmargin}{-0.5in}
\addtolength{\textheight}{1.0in}

\urlstyle{same}
\raggedbottom
\raggedright
\setlength{\tabcolsep}{0in}

%----------FORMATAÇÃO DE SEÇÕES----------
\titleformat{\section}{
  \vspace{-4pt}\scshape\raggedright\large
}{}{0em}{}[\color{black}\titlerule \vspace{-5pt}]

\pdfgentounicode=1

%----------COMANDOS CUSTOMIZADOS----------
\newcommand{\resumeItem}[1]{
  \item\small{#1 \vspace{-2pt}}
}

\newcommand{\resumeSubheading}[4]{
  \vspace{-2pt}\item
    \begin{tabular*}{0.97\textwidth}[t]{l@{\extracolsep{\fill}}r}
      \textbf{#1} & #2 \\
      \textit{\small#3} & \textit{\small #4} \\
    \end{tabular*}\vspace{-7pt}
}

% Linha de contexto em itálico abaixo do cabeçalho, sem marcador
\newcommand{\resumeDescription}[1]{%
  \begin{itemize}[leftmargin=0.15in, label={}]
    \item\small{\textit{#1}}
  \end{itemize}\vspace{-10pt}
}

\renewcommand\labelitemii{$\vcenter{\hbox{\tiny$\bullet$}}$}

\newcommand{\resumeSubHeadingListStart}{\begin{itemize}[leftmargin=0.15in, label={}]}
\newcommand{\resumeSubHeadingListEnd}{\end{itemize}}
\newcommand{\resumeItemListStart}{\begin{itemize}}
\newcommand{\resumeItemListEnd}{\end{itemize}\vspace{-5pt}}

%===========================================
\begin{document}

%----------CABEÇALHO----------
\begin{center}
  {\Huge \scshape Vinicius Cubas Brand} \\[4pt]
  \small
  Brasil (GMT$-$3) \,$|$\,
  \href{mailto:viniciuscb@gmail.com}{viniciuscb@gmail.com} \,$|$\,
  \href{https://linkedin.com/in/viniciusbrand}{linkedin.com/in/viniciusbrand} \,$|$\,
  \href{https://github.com/viniciuscb}{github.com/viniciuscb}
\end{center}

%----------RESUMO----------
\section{Resumo}
\begin{itemize}[leftmargin=0.15in, label={}]
  \small{\item{
    Engenheiro de Software Sênior Full-Stack e Mobile com 25~anos de experiência no desenvolvimento de
    aplicações web e mobile em SaaS, impacto social e sistemas embarcados. Proficiente em
    Ruby on Rails, React/TypeScript, Django/Python, React Native, etc, com histórico de conduzir
    funcionalidades complexas de ponta a ponta em equipes ágeis e remotas. Mestre em
    Informática. Baseado no Brasil (GMT$-$3).
  }}
\end{itemize}

%----------COMPETÊNCIAS----------
\section{Competências}
\begin{itemize}[leftmargin=0.15in, label={}]
  \small{\item{
    \textbf{Linguagens:} Ruby, Python, JavaScript/TypeScript, Java, PHP, SQL \\
    \textbf{Frameworks:} Ruby on Rails, Django, React, React Native, Redux, GraphQL (Apollo), Android (Java) \\
    \textbf{Bancos de Dados:} MySQL, PostgreSQL, Redis, Elasticsearch, CouchDB/PouchDB \\
    \textbf{Infraestrutura:} Linux, Docker, AWS (Lambda, EC2), Sidekiq, Git, Datadog, Airbrake \\
    \textbf{Ferramentas de IA:} Claude Code \\
    \textbf{Práticas:} Agile/Scrum, TDD, REST APIs, revisão de código, CI/CD \\
    \textbf{Idiomas:} Português (nativo), Inglês (avançado)
  }}
\end{itemize}

%----------EXPERIÊNCIA----------
\section{Experiência}
\resumeSubHeadingListStart

  %----- TapGoods -----
  \resumeSubheading
    {Engenheiro de Software Sênior}{Jan 2024 -- Presente}
    {TapGoods}{Remoto}
  \resumeDescription{Plataforma SaaS B2B multi-tenant atendendo centenas de empresas do setor de locação na América do Norte. Stack: Ruby on Rails, GraphQL, React, Redux, TypeScript, MySQL, Redis, Elasticsearch, Sidekiq, Docker.}
  \resumeItemListStart
    \resumeItem{Refatorou o módulo legado de roteamento e entregas como único programador --- colaborando com o líder técnico no relacionamento entre modelos e planejamento da migração, e coordenando com o suporte ao cliente para QA --- resolvendo um longo backlog de bugs reportados por clientes e melhorando significativamente a confiabilidade do roteamento.}
    \resumeItem{Implementou a funcionalidade de Agrupamento de Eventos de ponta a ponta (schema de banco de dados, API GraphQL, interface React e relatórios de faturamento), permitindo que operadores de locação rastreiem e faturem múltiplos pedidos em um único evento.}
    \resumeItem{Liderou uma iniciativa de padronização de impressos, migrando 3 documentos principais voltados ao cliente (faturas, resumos de locação, listas de separação) da renderização ad-hoc em React para um sistema unificado no backend Rails com templates HTML/CSS, incluindo refinamentos iterativos de estilo com a equipe de produto.}
    \resumeItem{Integrou uma biblioteca de leitura de código de barras mobile ao fluxo de trabalho de inventário serializado, permitindo que funcionários do armazém escaneiem itens com smartphones, eliminando a inserção manual de números de série.}
    \resumeItem{Implementou a funcionalidade Load Lists de ponta a ponta a partir de especificações de produto e designs no Figma: propôs o modelo de dados (diagrama ER), desenvolveu a interface full-stack e criou um sistema de impressão de etiquetas suportando tanto HTML/CSS quanto impressoras Zebra ZPL.}
    \resumeItem{Contribuiu em todo o backlog junto à equipe --- otimização de queries SQL, processamento assíncrono de imagens, tratamento de erros de API, triagem de bugs via Datadog e Airbrake, revisões de código em sprints ágeis diários --- e utiliza ativamente o Claude Code para desenvolvimento assistido por IA nas bases de código Rails e React.}
  \resumeItemListEnd

  %----- EITA -----
  \resumeSubheading
    {Desenvolvedor Full-Stack e Analista de Sistemas}{Dez 2013 -- Out 2024}
    {Cooperativa EITA}{Remoto}
  \resumeDescription{Cooperativa de tecnologia orientada por missão, desenvolvendo software para movimentos sociais, organizações de economia solidária e ONGs no Brasil.}
  \resumeItemListStart
    \resumeItem{Contribuiu com suporte a escrita LDAP ao projeto open-source Nextcloud, estendendo o core do Nextcloud para permitir a criação e o gerenciamento de usuários pela interface do Nextcloud quando as contas são mantidas em um diretório LDAP; a implementação serviu de base para o plugin oficial da comunidade \href{https://apps.nextcloud.com/apps/ldap_write_support}{ldap\_write\_support}.}
    \resumeItem{Atuou como principal desenvolvedor Ruby on Rails da equipe em múltiplos projetos de clientes, liderando decisões de arquitetura e entrega de ponta a ponta de sistemas baseados em RoR.}
    \resumeItem{Administrou a plataforma RIOS (\href{https://rios.eita.coop.br}{rios.eita.coop.br}), um serviço de nuvem SaaS multi-tenant para organizações sem fins lucrativos e movimentos sociais --- incluindo integração Nextcloud + Matrix.org (Synapse), funções AWS Lambda e gerenciamento de servidores Linux.}
    \resumeItem{Desenvolveu o Castanhadora, uma aplicação Android nativa (Java) que permite a coletores de castanhas na Amazônia brasileira registrar a produção e monitorar preços justos de venda.}
    \resumeItem{Desenvolveu o Lume, um sistema Python/Django que substituiu um fluxo de trabalho manual em planilhas para gerenciamento de propriedades agroecológicas, com visualização de dados e geração de relatórios.}
    \resumeItem{Colaborou diretamente com clientes não técnicos de movimentos comunitários para levantar requisitos, modelar dados e traduzir necessidades sociais em especificações técnicas em múltiplos projetos.}
  \resumeItemListEnd

  %----- Focus NFe frontend (part-time) -----
  \resumeSubheading
    {Engenheiro de Frontend (meio período)}{Jan 2013 -- Dez 2017}
    {Focus NFe}{Curitiba, Brasil}
  \resumeItemListStart
    \resumeItem{Desenvolveu a interface do usuário a partir de wireframes de design usando Ember.js, Bootstrap e CoffeeScript, integrando com uma API RESTful Ruby on Rails em uma equipe enxuta de 4 pessoas.}
  \resumeItemListEnd

  %----- RMI -----
  \resumeSubheading
    {Desenvolvedor Full-Stack}{Dez 2011 -- Out 2012}
    {RMI Informática Industrial}{Curitiba, Brasil}
  \resumeItemListStart
    \resumeItem{Como primeiro contratado da startup, desenvolveu componentes fundamentais do frontend usando Ruby on Rails, Backbone.js e Marionette.js.}
  \resumeItemListEnd

  %----- Dataprom -----
  \resumeSubheading
    {Desenvolvedor Full-Stack}{Fev 2008 -- Dez 2009}
    {Dataprom}{Curitiba, Brasil}
  \resumeItemListStart
    \resumeItem{Modernizou interfaces legadas e desenvolveu novas funcionalidades para um sistema de arrecadação automática de tarifas em larga escala (bilhetagem de transporte público) usando Java, Struts2 e Hibernate.}
  \resumeItemListEnd

  %----- Focus NFe full-stack -----
  \resumeSubheading
    {Desenvolvedor Full-Stack}{Set 2006 -- Jan 2008}
    {Focus NFe}{Curitiba, Brasil}
  \resumeItemListStart
    \resumeItem{Desenvolveu aplicações web full-stack para clientes das áreas de engenharia florestal e medicina usando PHP e Ruby on Rails, com consulta direta ao cliente.}
  \resumeItemListEnd

  %----- LACTEC -----
  \resumeSubheading
    {Desenvolvedor de Firmware}{Mar 2006 -- Set 2006}
    {LACTEC}{Curitiba, Brasil}
  \resumeItemListStart
    \resumeItem{Em parceria LACTEC--Siemens, desenvolveu firmware Linux embarcado em C para equipamento de telecomunicações com controle remoto.}
  \resumeItemListEnd

  %----- Início de carreira -----
  \resumeSubheading
    {Início de Carreira}{2000 -- 2005}{}{}
  \resumeItemListStart
    \resumeItem{Desenvolvedor full-stack na Thyamad Projects (PHP, customização eGroupWare); desenvolvedor de CMS e web na Editora Quantum (PHP, Oracle); Técnico em Eletrônica na Interatec; bolsista acadêmico na UFPR (ensino, pesquisa em processamento de imagens e planejamento em IA).}
  \resumeItemListEnd

\resumeSubHeadingListEnd

%----------FORMAÇÃO ACADÊMICA----------
\section{Formação Acadêmica}
\resumeSubHeadingListStart

  \resumeSubheading
    {Mestrado em Informática}{2007}
    {UFPR -- Universidade Federal do Paraná}{Curitiba, Brasil}
  \resumeItemListStart
    \resumeItem{Área: Processamento Digital de Imagens aplicado à Dermatologia}
  \resumeItemListEnd

  \resumeSubheading
    {Bacharelado em Ciência da Computação}{2004}
    {UFPR -- Universidade Federal do Paraná}{Curitiba, Brasil}

  \resumeSubheading
    {Técnico em Eletrônica}{2000}
    {UTFPR -- Centro Federal de Educação Tecnológica}{Curitiba, Brasil}

\resumeSubHeadingListEnd

%----------OUTRAS ATIVIDADES----------
\section{Outras Atividades}
\begin{itemize}[leftmargin=0.15in, label={}]
  \small{\item{
    Professor de yoga desde 2008 (segunda profissão). Ex-coordenador da CicloIguaçu, organização
    de ativismo cicloviário; co-organizador do III Fórum Mundial da Bicicleta (Curitiba, 2014,
    $\sim$3.000 participantes).
  }}
\end{itemize}

\vspace{10pt}
\begin{center}
  \small\textit{Última atualização: Abril de 2026}
\end{center}

\end{document}
